% Course additions; preamble.tex follows the Robotics template.
\usepackage{tabularx}
\usepackage{needspace}
\IfFontExistsTF{Source Han Sans}
  {\setCJKsansfont{Source Han Sans}}
  {\setCJKsansfont{FandolHei-Regular.otf}[BoldFont=FandolHei-Bold.otf]}
\pgfplotsset{compat=1.18}
\hypersetup{
  pdftitle={机器人传感技术：第2章信号分析与描述复习册},
  pdfauthor={魏崃},
  pdfsubject={信号分类、傅里叶分析、典型例题与复习自测}
}
\newcommand{\iu}{\mathrm{j}}
\newcommand{\diff}{\mathrm{d}}
\DeclareMathOperator{\Sa}{Sa}
\DeclareMathOperator{\sgn}{sgn}
\DeclareMathOperator{\PV}{PV}
\newcommand{\keypoint}[1]{\textcolor{myr}{\textbf{#1}}}
% The examples and theorems in this chapter are short enough to keep intact.
\renewcommand{\ex}[2]{\begin{Example}[breakable=false]{#1}{}#2\end{Example}}
\renewcommand{\thm}[2]{\begin{Theorem}[breakable=false]{#1}{}#2\end{Theorem}}
\newtcolorbox{focusbox}[1]{enhanced,breakable,colback=whusakura!10,
  colframe=whusakura!85!black,boxrule=.6pt,sharp corners,
  title={#1},fonttitle=\bfseries,coltitle=black,colbacktitle=whusakura!45}
